% Page 23 — Alban – XIIIe siècle

\begin{center}
  {\Large\itshape Alban – XIII\textsuperscript{ème} siècle}
\end{center}
\vspace{1.2em}

\noindent
\lettrine{B}{ien sûr}, la légende « enrichie » de Notre Dame d’Alban est déjà un témoignage de la ferveur populaire… et il y a bien de poésie et d’humour en ces siècles dits « de fer ». Mais nous aimerions tant connaître les lieux et les hommes de ce temps-là. Nous n’avons pour cela que quelques pierres, quelques bribes d’archives. Deux chapiteaux romans – dont l’un avec bélier – retrouvés et replacés en remploi dans l’encadrement du portail actuel de l’église ne suffisent pas à nous silhouetter l’humble sanctuaire qui se dressait entre les murs de notre petite cité médiévale. De maigres documents n’apaiseront pas notre curiosité pour suivre le passage d’une civilisation à une autre. Mais déjà s’ébauche la disparition du servage.

\medskip

\noindent
\lettrine{E}{n dehors} de la noblesse, il y a quantité d’hommes libres, mais aussi une non moindre quantité d’individus de condition un peu imprécise entre la liberté et la servitude. Lisons quelques pièces d’archives en prenant garde de nous tromper sur leur portée. Le 23 décembre 1248, Isarn de Paulin vend à l’évêque d’Albi la sixième partie des revenus du château de Nougaret, ses dépendances, « serfs compris ». Le 12 juillet 1249, les frères Vassal font vente au même prélat, en des termes qui font rêver : « Je vous vends plus spécialement Laurent et ses enfants, et la moitié de Bernard et de ses enfants… » et en texte occitan :

\medskip

\noindent
\lettrine{E}{n langue} d’Oc : « vindé mai ochéna (8°) art dél castèl de Nogarét et la ochena part del moli et del fourri, et la ochena part de las semorias dels homes et de las femnas dél dig castèl… » Que cela ne nous passe pas dresser les cheveux sur la tête. Nous devinons qu’il n’est pas question entre acheteur et vendeur de partager en deux pièces Laurent, ou Bernard, les hommes et les femmes du coin, mais de se répartir les revenus fournis par eux.
